\chapter{后训练方法}
\label{lab:41_post_training}

\noindent\textbf{配套文件：}\path{notebook/41_post_training.ipynb}\par
\noindent\textbf{教材关联：}\bookxrefshort{ch:supervised-finetuning}、\bookxrefshort{ch:reinforcement-learning}、\bookxrefshort{ch:preference-alignment}、\bookxrefshort{ch:verifiable-reasoning}。

\section*{实验目标}
分别验证监督区域、偏好概率比和组内优势的数值对象。

\section*{准备及定位}
先阅读本册实验总则，确认Notebook中的依赖与资产单元。原理计算、库接口迁移和真实模型训练可能具有不同资源要求，应分别记录设备、版本及运行范围。

源码定位可从下列函数开始；应结合前置数据与配置单元阅读。
\begin{itemize}
\item \texttt{my\_format\_prompt}
\item \texttt{my\_encode\_sft\_record}
\item \texttt{my\_padding\_report}
\end{itemize}

\section*{操作及观察}
\begin{enumerate}
\item 手工标注一条指令回答的有效标签，检查动态padding及packing后的监督区域。
\item 运行SFT基线并记录有效词元归一化方式。
\item 对固定偏好样本计算DPO目标，区分策略与参考模型的序列对数概率。
\item 对固定奖励组核对GRPO优势、裁剪及零方差处理，再查看TRL对象映射。
\end{enumerate}

\section*{结果判读及提交}
提交监督掩码、偏好计算和组内优势表。固定奖励数值验证只证明目标计算关系，不等于已经完成在线强化学习训练。

提交时附上所执行单元范围、环境与制品身份、原始结果及失败说明。将未运行项目明确列出，不能用Notebook中已有输出替代本次执行证据。

相关方法的原始定义与适用前提见\cite{rafailov2023dpo,shao2024deepseekmath}；本实验的运行结果仍须按实际配置单独验证。
